\documentclass[11pt,a4paper]{article}

% --- PAQUETS DE BASE ---
\usepackage[utf8]{inputenc}
\usepackage[T1]{fontenc}
\usepackage[french]{babel}
\usepackage{amsmath, amsfonts, amssymb}
\usepackage{geometry}
\geometry{top=2cm, bottom=2.5cm, left=2cm, right=2cm}

% --- DESIGN & COULEURS ---
\usepackage{xcolor}
\usepackage{colortbl}
\definecolor{BrandPrimary}{HTML}{0D9488}
\definecolor{BrandDark}{HTML}{111827}

\usepackage{tcolorbox}
\tcbuselibrary{skins, breakable}

\usepackage{titlesec}
\titleformat{\section}{\color{BrandPrimary}\normalfont\Large\bfseries}{\thesection}{1em}{}[{\titlerule[0.8pt]}]
\titleformat{\subsection}{\color{BrandDark}\normalfont\large\bfseries}{\thesubsection}{1em}{}

% --- POLICE ---
\usepackage{helvet}
\renewcommand{\familydefault}{\sfdefault}

\newcommand{\makebrandtitle}{
    \begin{tcolorbox}[enhanced, colback=BrandDark, colframe=BrandPrimary, arc=10pt, boxrule=2pt, halign=center]
        \vspace{0.3cm}
        {\Huge \bfseries \color{white} ZENITH LEARN} \\
        \vspace{0.2cm}
        {\Large \color{BrandPrimary} IA Éthique : Modélisation TWIST 10} \\
        \vspace{0.3cm}
    \end{tcolorbox}
    \vspace{0.5cm}
}

\begin{document}

\makebrandtitle

\section{Introduction et Problématique}
Le \textbf{TWIST 10} aborde la question cruciale de la fiabilité du signal : \textit{``Les données d'action terrain sont rares et parfois auto-déclarées.''} 

Dans une plateforme orientée vers l'impact, le risque de fraude est immense dès lors qu'un financement est en jeu (TWIST 07). Une donnée auto-déclarée (ex: ``J'ai contacté 10 partenaires'') a une valeur de preuve largement inférieure à une donnée vérifiée par le système (ex: ``L'apprenant a complété son module''). Se fier aveuglément à l'auto-déclaration créerait une \textbf{dépendance destructrice} : le système récompenserait les meilleurs menteurs plutôt que les meilleurs entrepreneurs.

\section{Objectif Apparent vs Réalité Algorithmique}
\begin{itemize}
    \item \textbf{Objectif Apparent :} Utiliser toutes les données disponibles pour évaluer le potentiel.
    \item \textbf{Réalité Algorithmique :} Favorise ceux qui déclarent le plus d'actions, quel que soit le degré de véracité, créant un biais de "volontarisme autodéclaré".
\end{itemize}

\section{Modélisation de la Confiance ($CR$)}
L'IA doit désormais douter. Nous introduisons l'indice de confiance de la donnée ($CR$) comme multiplicateur de fiabilité.

\subsection{Formule du Score de Fidélité ($SPI^{final}$)}
\begin{equation}
SPI^{final}(u) = \sum_{j \in \text{Features}} (F_j(u) \cdot w_j \cdot CR_j)
\end{equation}

\subsection{Genèse de la formule : pourquoi la pondération de confiance ?}
La formule exprime que la valeur d'une information ($F_j$) dépend de la fiabilité de son canal de transmission ($CR_j$). 

Plutôt que d'exclure les données auto-déclarées (ce qui serait dommageable car elles contiennent une part de vérité), nous appliquons un \textbf{coefficient de prudence}. C'est une approche bayésienne : nous avons un "a priori" de confiance variable selon la source de l'information. Un succès "prouvé" vaut mathématiquement plus qu'un succès "raconté".

\subsection{Détail et justification des termes}
\begin{itemize}
    \item \textbf{$F_j(u) \cdot w_j$ :} Le score de base pondéré par l'influence du segment (T05).
    
    \item \textbf{$CR_j$ (Confidence Rating) :} Le coefficient de décote appliqué à la source de la donnée $j$.
    \begin{itemize}
        \item \textbf{Verified (Provenance Système)} : $CR = 1.0$. Confiance totale (ex: temps passé, logs de connexion).
        \item \textbf{Self-Reported (Déclaré apprenant)} : $CR = 0.7$. Nous appliquons une \textbf{prudence de 30\%}. Pour que le score auto-déclaré d'un apprenant atteigne la valeur d'un score système, il doit en faire 40\% plus.
        \item \textbf{Imputed (Estimé/Rare)} : $CR = 0.35$. Pour les 80\% d'explorateurs sans historique, l'IA estime des valeurs par voisinage mais leur accorde une confiance très faible.
    \end{itemize}

    \item \textbf{Effet sur la fraude :} Un apprenant qui gonfle ses déclarations verra son score plafonné par son $CR_j$. Pour entrer dans le Top 100, il devra nécessairement valider des signaux de confiance forte (système, audit externe).
\end{itemize}

\section{Structure du Dataset (Audit de Fiabilité)}
Le dataset intègre désormais l'origine de chaque point de donnée.
\begin{table}[h!]
\centering
\begin{tabular}{|l|l|p{7cm}|}
\hline
\rowcolor{BrandDark} \color{white} Colonne & \color{white} Type & \color{white} Rôle T10 \\ \hline
\texttt{is\_self\_reported} & Bool & Déclencheur du coefficient $CR = 0.7$. \\ \hline
\rowcolor{BrandPrimary!20} \texttt{action\_integrity} & Float & Preuve de fiabilité pour le bailleur. \\ \hline
\texttt{source\_audit\_log} & String & Traçabilité du canal de donnée. \\ \hline
\end{tabular}
\end{table}

\section{Conclusion}
Le TWIST 10 ferme la boucle de l'IA de Zenith Learn. En apprenant à l'algorithme à \textbf{douter et à vérifier}, nous protégeons l'intégrité du système de financement. Nous ne finance pas seulement le potentiel, nous finançons la \textbf{Vérité Entrepreneuriale}. Zenith Learn devient ainsi une plateforme de confiance, capable de garantir au bailleur que chaque apprenant sélectionné l'a été sur la base de faits tangibles et pondérés avec sagesse.

\end{document}
